% ECONOMICS_PROBLEM: {"id": "H11-328", "title": "Building effective state capacity", "jel_code": "H11", "source_id": "OP-0032"}
% !TeX program = xelatex
\documentclass[11pt,a4paper]{article}
\usepackage[margin=23mm]{geometry}
\usepackage{fontspec}
\setmainfont{Latin Modern Roman}
\usepackage{amsmath,amssymb,mathtools}
\usepackage{xcolor,enumitem,booktabs,longtable,array,xurl,hyperref}
\definecolor{muted}{HTML}{53636C}
\hypersetup{colorlinks=true,urlcolor=blue,linkcolor=blue}
\setlength{\parindent}{0pt}
\setlength{\parskip}{5pt}
\newcommand{\R}{\mathbb R}
\newcommand{\E}{\mathbb E}
\newcommand{\N}{\mathbb N}
\newcommand{\ind}{\mathbf 1}
\newcommand{\Var}{\operatorname{Var}}
\newcommand{\Cov}{\operatorname{Cov}}
\newcommand{\supp}{\operatorname{supp}}
\DeclareMathOperator*{\argmin}{arg\,min}
\DeclareMathOperator*{\argmax}{arg\,max}
\newcommand{\entrymeta}[3]{{\sffamily\small\color{muted}\textbf{\texttt{#1}}\quad|\quad #2\par #3\par}}
\newcommand{\Problemref}[1]{\texttt{#1}}
\newcommand{\JELref}[2]{\texttt{#2} (JEL \texttt{#1})}
\begin{document}
\section*{Building effective state capacity}
\textbf{Problem ID:} \texttt{H11-328}\quad \textbf{JEL:} \texttt{H11}\par
% BEGIN ECONOMICS STATEMENT
\label{problem:OP-0032}
\label{atlas:prob:D2}
\noindent\textbf{Original atlas ID:} D2 (entry 32). \textbf{Classification:} Editorial politically constrained capacity-investment problem.

\noindent\textbf{Original atlas question:} Which investments create durable fiscal, legal and administrative capability without empowering predation?

\subsection*{Definitions and assumptions}
Consider a dynamic political economy with state $S_t=(k_t,z_t,p_t)$, where $k_t=(k_t^f,k_t^l,k_t^a)$ records fiscal, legal, and administrative capabilities, $z_t$ records economic conditions, and $p_t$ records political control. Capability coordinates require operational measures, such as audit accuracy or case-resolution capacity, rather than tax receipts alone. A feasible investment rule $u$ maps observed histories into training, records, courts, and technology investments. Model $M$ specifies the transition kernel, private incentives, political turnover, information, resource constraints, and an equilibrium selection rule. Let $C_t^M(u)$ be real public-service benefit net of investment costs and $R_t^M(u)\geq0$ a prespecified measure of predation or coercive loss.
\subsection*{Formal research question}
For fixed $H$, $\beta\in(0,1)$, welfare weights, and tolerance $\bar R\geq0$, characterize
\[
 \sup_{u\in\mathcal U_M}
     \mathbb E_M\sum_{t=0}^H\beta^t C_t^M(u)
 \quad\text{subject to}\quad
     \mathbb E_M\sum_{t=0}^H\beta^t R_t^M(u)\leq\bar R,
\]
where $\mathcal U_M$ contains budget-feasible, politically implementable rules under the announced equilibrium concept. Identify the policy response of capability coordinates and welfare across all models consistent with administrative and experimental data. Determine which investments yield persistent gains after support ends or control changes, and which merely improve incumbent extraction or measured collection.
\subsection*{Interpretation and scope}
State capacity is not a scalar synonym for benevolent governance. The same records and enforcement technology can improve service provision or facilitate coercion; the objective and constraint make that tradeoff explicit. Observed historical persistence does not identify the effect of introducing a capability technology today. Randomizing an administrative tool identifies an assignment effect at the randomized level but can change politician, bureaucrat, and taxpayer behavior outside that unit. Such spillovers and organizational incentives need measurement or structural assumptions. Fiscal capacity is distinct from a tax-rate increase, and higher revenue may reflect a growing base rather than administrative improvement. Durability should be tested under a defined horizon and turnover process, not assumed from a single follow-up.
\subsection*{Source audit and status}
[S1] models origins of capacity, [S2] links capacity and conflict, and [S3] studies network interactions. The atlas's linked NBER Working Paper 15088 is dated 2009; the cited Econometrica article is 2010. The constrained policy problem above is an editorial synthesis, not a theorem posed by those papers. Older theoretical and empirical results do not certify general unresolved status in 2026.

\subsection*{References}
\begin{enumerate}[label=\textnormal{[S\arabic*]},leftmargin=*]
\item Timothy Besley and Torsten Persson (2009). \emph{The Origins of State Capacity: Property Rights, Taxation, and Politics}. American Economic Review 99(4), 1218--1244. \url{https://doi.org/10.1257/aer.99.4.1218}. \textit{Locator:} Primary publisher, working-paper, or author record: bibliographic information and abstract; full-text theorem verification is not claimed. \textit{Role:} Model of capacity investment and political incentives.
\item Timothy Besley and Torsten Persson (2010). \emph{State Capacity, Conflict, and Development}. Econometrica 78(1). \url{https://doi.org/10.3982/ECTA8073}. \textit{Locator:} Primary publisher, working-paper, or author record: bibliographic information and abstract; full-text theorem verification is not claimed. \textit{Role:} Joint fiscal/legal capacity and conflict framework; original link to WP15088 denotes the earlier 2009 version.
\item Daron Acemoglu, Camilo García-Jimeno, and James A. Robinson (2015). \emph{State Capacity and Economic Development: A Network Approach}. American Economic Review 105(8). \url{https://doi.org/10.1257/aer.20140044}. \textit{Locator:} Primary publisher, working-paper, or author record: bibliographic information and abstract; full-text theorem verification is not claimed. \textit{Role:} Network model and evidence on local state-capacity interactions.
\end{enumerate}

\subsection*{Common conventions: Source mathematical conventions}

Unless an entry states otherwise, agent and firm sets are finite. State and action spaces are measurable, strategies and outcome maps are measurable, probability laws are specified, and expectations exist for the targets under discussion. Infinite-horizon discount factors lie in $(0,1)$ and discounted objectives are finite. Norms, tolerances and complexity input sizes must be specified for computational comparisons. Constants or primitives that appear locally are inputs, not empirical facts asserted by this catalogue.

\textbf{Identification.} A statistical model $m\in\mathcal M$ implies an observable law $P_m$ and target $\theta(m)$. At law $P$, its identified set is
\[
\Theta(P;\mathcal M)=\{\theta(m):m\in\mathcal M,\ P_m=P\}.
\]
Point identification holds when this set is a singleton, or equivalently when $P_{m_1}=P_{m_2}$ implies $\theta(m_1)=\theta(m_2)$ for all compatible models. Sharp bounds mean bounds attained, or approached when the boundary is not attained, by compatible models. Writing this set is a definition, not a solution: the substantive task is to establish informative restrictions and derive or estimate the set. An unrestricted-confounding model may give only trivial bounds.

\textbf{Inference.} Identification, estimability and valid uncertainty quantification are separate questions. Fix $\alpha\in(0,1)$. A confidence procedure $C_n$ over a declared law class $\mathcal P$ has asymptotic uniform coverage at level $1-\alpha$ when
\[
\liminf_{n\to\infty}\inf_{P\in\mathcal P}\Pr_P\{\theta(P)\in C_n\}\geq1-\alpha.
\]
For set-identified targets the coverage event must be defined separately, for example coverage of the full identified set. If an entry uses identification-indexed models instead of uniquely indexed laws, the same coverage convention is applied over those models. Pointwise limits do not establish uniform validity, and asymptotic validity does not guarantee finite-sample precision. Sample sizes, dependence structures and weak-identification sequences are part of each application.

\textbf{Counterfactuals.} $Y_i(a)$ is a potential outcome under a specified intervention, including its timing, implementation and versions. It is not $Y_i$ conditional on voluntarily choosing $a$. A full intervention vector permits interference; shorthand scalar treatment notation does not silently impose no interference. Assignment, selection, attrition, anticipation and measurement must be specified. A claim about an instrument's local effect is not automatically a population policy effect.

\textbf{Economic equilibrium.} A counterfactual model includes best responses, constraints, feasibility, market clearing, state transitions and expectations consistent with the stated information. When several equilibria exist, either a declared selector is an input or the target is set-valued. Comparisons across policy regimes must use the stated selection convention. Reduced-form relationships are not presumed invariant after an equilibrium-changing intervention.

\textbf{Optimization and welfare.} Optimization is over the stated feasible set. If attainment is not established, $\sup$ or $\inf$ and $\varepsilon$-optimal choices replace an asserted maximizer or minimizer. For example, with value $V=\sup_{a\in\mathcal A}W(a)<\infty$, an $\varepsilon$-optimal set is $\{a\in\mathcal A:W(a)\geq V-\varepsilon\}$ for $\varepsilon>0$. Benefits, costs, risk and health or environmental outcomes can be aggregated only using declared units and conversion weights. Interpersonal utility weights, the treatment of fiscal transfers and foreign profits, discounting, and the behavioral welfare criterion are explicit inputs. A comparative-static derivative additionally requires existence and appropriate differentiability of the selected counterfactual.

\textbf{Sources.} Local labels [S1], [S2], and so forth restart in every question. Locators identify the relevant abstract, model, section or result at the level actually checked; a bibliographic or abstract-level verification is not represented as inspection of an entire proof. Working papers are distinguished from later publications and changing versions. Source summaries are paraphrased. All URL checks and editorial formulations refer to the audit date stated above.
% END ECONOMICS STATEMENT
\end{document}
